\chapter{Volatilité locale}\label{ch:dv:local-volatility}

Deux desks valorisent le même put à départ différé sur un indice~: le strike est fixé à 95~\% du niveau
qu'atteindra l'indice dans six mois, et l'échéance tombe six mois plus tard.
Les deux utilisent des modèles qui reproduisent chaque vanille cotée au tick près.
Le premier modèle affirme que, dans six mois, le smile à six mois sera à peine plus de la moitié
aussi pentu qu'aujourd'hui, et il valorise le put à bas prix~;
le second affirme que le smile ressemblera à ce qu'il est aujourd'hui, et il facture davantage.
Les deux s'accordent avec le marché d'aujourd'hui~; ils divergent sur celui de demain, et une option à départ
différé est un pari sur le smile de demain. Le premier modèle est le \omterm{def:dv:local-volatility:model}{modèle de volatilité locale}~:
la seule diffusion dont la volatilité est une fonction du temps et du spot, et qui reproduit le prix de chaque option européenne.
Ce chapitre la déduit de la surface (\omterm{thm:dv:local-volatility:dupire}{formule de Dupire}), montre pourquoi elle constitue le
point de départ naturel de tout modèle de smile, la calibre et la vérifie par simulation, et
mesure ce qu'elle annonce pour le futur~: le \omterm{def:dv:local-volatility:forward}{smile forward} et la manière dont le smile se
déplace lorsque le spot se déplace.

\section{Formule de Dupire}

\begin{definition}[Modèle de volatilité locale]\label{def:dv:local-volatility:model}
Dans un \emph{modèle de volatilité locale}\index{modèle de volatilité locale} le sous-jacent suit
\[
dS_t=(r-q)S_t\,dt+\sigma_{\mathrm{loc}}(t,S_t)S_t\,dW_t
\]
sous la probabilité risque-neutre, avec une fonction déterministe $\sigma_{\mathrm{loc}}$, sa
\emph{volatilité locale}\index{volatilité locale}. Le modèle est complet~: la seule source de
risque est le spot, et tout \omterm{def:dv:no-arbitrage-and-the-fundamental-theorems:claim}{actif contingent} est répliqué avec le sous-jacent.
\end{definition}

Le \omterm{def:dv:black-scholes-three-ways:model}{modèle de Black--Scholes} correspond au cas d'un $\sigma_{\mathrm{loc}}$ constant. La question est
de savoir si une fonction reproduit toute une surface de prix d'options, et laquelle. La réponse de Dupire
est qu'une seule le fait, et qu'elle se lit directement sur les prix.

\begin{theorem}[Formule de Dupire]\label{thm:dv:local-volatility:dupire}
Si les prix des calls $C(K,T)$ sont réguliers en $(K,T)$ et engendrés par un \omterm{def:dv:local-volatility:model}{modèle de volatilité locale}, alors
\[
\sigma_{\mathrm{loc}}^2(T,K)=\frac{\partial_TC+(r-q)K\,\partial_KC+qC}{\tfrac12K^2\,\partial_{KK}C},
\]
la \emph{formule de Dupire}\index{formule de Dupire}. Réciproquement, toute surface régulière de prix de calls
sans arbitrage détermine une \omterm{def:dv:local-volatility:model}{volatilité locale} qui la reproduit.
\end{theorem}
\begin{proof}[Démonstration partielle (taux et dividendes nuls)]
La densité $q(T,s)$ de $S_T$ est solution de l'équation de Kolmogorov progressive de la diffusion (livre~4,
chapitre~4)~: $\partial_Tq=\tfrac12\partial_{ss}\bigl(\sigma_{\mathrm{loc}}^2s^2q\bigr)$.
Écrivons le call sous la forme
\[
C(K,T)=\int_K^\infty(s-K)\,q(T,s)\,ds,
\]
on dérive en $T$, on substitue et on intègre par parties deux fois~:
$\partial_TC=\tfrac12\sigma_{\mathrm{loc}}^2(T,K)K^2q(T,K)$. Avec $q=\partial_{KK}C$
(Breeden--Litzenberger, chapitre~7), on obtient la formule.
\end{proof}

Le numérateur est le prix d'un calendar spread, le dénominateur celui d'un butterfly~: la variance locale est
le rapport des deux \omterm{def:dv:implied-volatility-and-its-surface:static}{arbitrages statiques} du chapitre~7, et elle est positive exactement lorsque les deux le sont.

\section{Volatilité locale à partir de la volatilité implicite}

Les prix sont des entrées inadaptées pour les dérivés~: ils couvrent plusieurs ordres de grandeur selon les strikes, et une
dérivée seconde de prix bruités n'est que du bruit. La même formule, exprimée dans les coordonnées du chapitre~7,
se comporte beaucoup mieux.

\begin{proposition}[Dupire en variance implicite totale]\label{prop:dv:local-volatility:w}
Avec $k=\ln(K/F_{0,T})$ et $w(k,T)=\sigma_{\mathrm{imp}}^2(k,T)T$,
\[
\sigma_{\mathrm{loc}}^2(T,\,F_{0,T}e^k)=\frac{\partial_Tw}{g(k)},\qquad
g(k)=\Bigl(1-\frac{k\,\partial_kw}{2w}\Bigr)^2-\frac{(\partial_kw)^2}{4}\Bigl(\frac1w+\frac14\Bigr)
+\frac{\partial_{kk}w}2 ,
\]
où $g$ est le facteur de butterfly du chapitre~7.
\end{proposition}
\admitted

\begin{proposition}[Le skew local vaut deux fois le skew implicite]\label{prop:dv:local-volatility:twice}
Pour des échéances courtes et un skew faible, la volatilité implicite au strike $K$ est proche de la moyenne de la
\omterm{def:dv:local-volatility:model}{volatilité locale} entre le spot et $K$~; en particulier, la pente de la \omterm{def:dv:local-volatility:model}{volatilité locale} en log-strike à la
monnaie vaut le double de la pente de la volatilité implicite.
\end{proposition}
\begin{proof}[Esquisse]
Pour une échéance courte, le chemin de $S$ à $K$ est proche d'une droite en $\ln S$, et la variance implicite est la
moyenne de la variance locale le long de ce chemin. Si $\sigma_{\mathrm{loc}}$ est linéaire en $\ln K$ de pente $2a$, sa
moyenne de $\ln S$ à $\ln K$ est linéaire de pente $a$.
\end{proof}

Sur la surface du chapitre (une surface SSVI du type de celle calibrée au chapitre~8, avec taux et
dividendes nuls), la volatilité implicite à trois mois vaut 16.4~\% à la monnaie et 21.1~\% pour $k=-0.1$~; la
\omterm{def:dv:local-volatility:model}{volatilité locale} vaut 17.7~\% et 27.4~\%. Les pentes à la monnaie valent $-0.47$ et $-0.92$ par unité de
\omterm{def:dv:implied-volatility-and-its-surface:surface}{log-moneyness}, soit un rapport de 1.98 (\cref{fig:dv:local-volatility:lv}).

\begin{omfigure}
\begin{tikzpicture}
\begin{axis}[width=10.5cm,height=5cm,
  xlabel={log-moneyness $k$},ylabel={volatilité (\%)},
  tick label style={font=\small},label style={font=\small},
  xmin=-0.4,xmax=0.3,ymin=10,ymax=46,grid=major,grid style={gray!25},
  legend columns=2,legend style={font=\footnotesize,at={(0.5,-0.3)},anchor=north,draw=none,fill=none,/tikz/every even column/.append style={column sep=8pt}},
  table/col sep=comma]
\addplot[omDef,thick] table[x=k,y=imp3m]{figdata/derivatives/09-local-volatility/local_vs_implied.csv};
\addplot[omDef,very thick,dashed] table[x=k,y=loc3m]{figdata/derivatives/09-local-volatility/local_vs_implied.csv};
\addplot[omThm,thick] table[x=k,y=imp1y]{figdata/derivatives/09-local-volatility/local_vs_implied.csv};
\addplot[omThm,very thick,dashed] table[x=k,y=loc1y]{figdata/derivatives/09-local-volatility/local_vs_implied.csv};
\legend{volatilité implicite 3 mois,volatilité locale 3 mois,volatilité implicite 1 an,volatilité locale 1 an}
\end{axis}
\end{tikzpicture}

\omcaption{Volatilité implicite (trait plein) et \omterm{def:dv:local-volatility:model}{volatilité locale} de Dupire (tirets) sur la
surface du chapitre, à trois mois et à un an. Au voisinage de la monnaie, la \omterm{def:dv:local-volatility:model}{volatilité locale} est
deux fois plus pentue que l'implicite~; sur les ailes, elle épouse la courbure du smile. Données~:
code du chapitre.}\label{fig:dv:local-volatility:lv}
\end{omfigure}

\section{Calibration en pratique}

La \omterm{thm:dv:local-volatility:dupire}{formule de Dupire} demande les dérivées d'une surface, et non des cotations. Trois règles en découlent.

\begin{method}[Construction d'une surface de volatilité locale]\label{met:dv:local-volatility:build}
\begin{enumerate}
\item Ajuster d'abord une surface implicite lisse et sans arbitrage (SVI ou SSVI, chapitre~8), après avoir retiré les événements
et appliqué le \omterm{def:dv:parametrising-the-surface:business}{temps ouvré}~; la formule renvoie alors une variance locale positive et finie.
\item Évaluer la \omterm{def:dv:local-volatility:model}{volatilité locale} sur une grille en temps et en \omterm{def:dv:implied-volatility-and-its-surface:surface}{log-moneyness} par rapport au forward de la
date de calibration, puis l'interpoler~; au-delà de la plage ajustée, prolonger à plat.
\item Vérifier le résultat en valorisant les vanilles avec le \omterm{def:dv:local-volatility:model}{modèle de volatilité locale} (par grille ou par Monte-Carlo)
et en les comparant à la surface~; des écarts au-delà du bruit numérique signifient que la grille ou l'extrapolation
est trop grossière.
\end{enumerate}
\end{method}

La \Cref{fig:dv:local-volatility:reprice} montre la vérification~: des vanilles à six mois valorisées par simulation de 200\,000
trajectoires du \omterm{def:dv:local-volatility:model}{modèle de volatilité locale}, avec des pas quotidiens, restituent les volatilités implicites de la surface à
moins de 0.11 point de volatilité près. Les erreurs ont toutes le même signe~: un biais d'environ 0.08 point provient du
pas de temps quotidien~: avec quatre pas par jour, elles tombent sous 0.06 point dans un sens comme dans l'autre, tandis
qu'une grille quatre fois plus fine ne les modifie que de 0.01 à 0.02.

\begin{omfigure}
\begin{tikzpicture}
\begin{axis}[width=10.5cm,height=5cm,
  xlabel={strike},ylabel={volatilité implicite à 6 mois (\%)},
  tick label style={font=\small},label style={font=\small},
  xmin=72,xmax=128,ymin=12,ymax=30,grid=major,grid style={gray!25},
  legend columns=2,legend style={font=\footnotesize,at={(0.5,-0.3)},anchor=north,draw=none,fill=none,/tikz/every even column/.append style={column sep=8pt}},
  table/col sep=comma]
\addplot[omDef,very thick] table[x=strike,y=vol]{figdata/derivatives/09-local-volatility/surface6m.csv};
\addplot[omThm,only marks,mark=*,mark size=1.6pt,error bars/.cd,y dir=both,y explicit]
  table[x=strike,y=mc,y error expr=\thisrow{hi}-\thisrow{mc}]{figdata/derivatives/09-local-volatility/reprice.csv};
\legend{surface,volatilité locale Monte-Carlo ($\pm2$ erreurs type)}
\end{axis}
\end{tikzpicture}

\omcaption{La vérification de la \cref{met:dv:local-volatility:build}~: des options à six mois valorisées par simulation
du \omterm{def:dv:local-volatility:model}{modèle de volatilité locale} (200\,000 trajectoires, pas quotidiens) face à la surface implicite à laquelle elles ont été
calibrées. Données~: le tutoriel.}\label{fig:dv:local-volatility:reprice}
\end{omfigure}

La \omterm{def:dv:local-volatility:model}{volatilité locale} est aussi une interprétation. Tout modèle dans lequel le spot a une volatilité stochastique
a les mêmes lois unidimensionnelles qu'un certain \omterm{def:dv:local-volatility:model}{modèle de volatilité locale}.

\begin{definition}[Projection markovienne]\label{def:dv:local-volatility:projection}
La \emph{projection markovienne}\index{projection markovienne} d'un processus d'Itô $dS_t=\mu_tS_t\,dt+
\sigma_tS_t\,dW_t$, avec $\sigma_t$ aléatoire, est le \omterm{def:dv:local-volatility:model}{modèle de volatilité locale} tel que
$\sigma_{\mathrm{loc}}^2(t,K)=\E\bigl[\sigma_t^2\mid S_t=K\bigr]$~; d'après le théorème de Gyöngy, $S_t$ a la
même loi dans les deux modèles à chaque date.
\end{definition}

Tout modèle de smile possède donc une \omterm{def:dv:local-volatility:model}{volatilité locale}~: l'espérance conditionnelle de sa variance instantanée
sachant le spot. C'est ainsi que se calibrent les modèles à \omterm{def:dv:local-volatility:model}{volatilité locale} stochastique (chapitre~20) et la
corrélation locale (chapitre~17), et pourquoi ajuster les vanilles du jour ne dit rien de la dynamique d'un modèle~:
tout modèle qui les ajuste possède la même \omterm{def:dv:local-volatility:projection}{projection markovienne}.

\section{Ce que prédit la volatilité locale~: dynamique et smile forward}

Un \omterm{def:dv:local-volatility:model}{modèle de volatilité locale} est ajusté au smile du jour~; tout ce qu'il dit de demain découle de
cet ajustement. Deux conséquences importent.

\begin{proposition}[Dynamique du smile sous volatilité locale]\label{prop:dv:local-volatility:dynamics}
Sous \omterm{def:dv:local-volatility:model}{volatilité locale}, lorsque le spot baisse, la volatilité implicite de chaque strike fixe augmente (le smile
se déplace vers les strikes plus élevés), et la volatilité à la monnaie se déplace, pour de petits mouvements et de faibles skews, deux fois
aussi vite que sous \omterm{def:dv:implied-volatility-and-its-surface:sticky}{sticky strike}~: $\partial\sigma_{\mathrm{ATM}}/\partial\ln S=2\,\partial\sigma_{\mathrm{imp}}
/\partial\ln K$ à la monnaie.
\end{proposition}
\begin{proof}
Par la \cref{prop:dv:local-volatility:twice}, la volatilité implicite à la monnaie est proche de la \omterm{def:dv:local-volatility:model}{volatilité locale}
au spot, dont la pente en $\ln S$ vaut deux fois le skew implicite~; la volatilité implicite d'un strike fixe
est la moyenne de la \omterm{def:dv:local-volatility:model}{volatilité locale} entre le nouveau spot et le strike, qui augmente lorsque le spot se déplace vers
une zone de \omterm{def:dv:local-volatility:model}{volatilité locale} plus élevée.
\end{proof}

Sur le smile à trois mois du chapitre, l'effet dépasse celui de la règle~: une baisse de 1~\% de l'indice fait monter
la volatilité à la monnaie 2.7 fois plus que sous \omterm{def:dv:implied-volatility-and-its-surface:sticky}{sticky strike}, et une baisse de 5~\% 2.9 fois plus, de 16.4~\% à
23.4~\% contre 18.8~\% (\cref{fig:dv:local-volatility:dynamics}). La règle vaut pour les échéances courtes et les faibles
skews~; le skew à trois mois du chapitre n'est ni faible ni très court. Hagan et ses coauteurs ont montré que sur
les marchés de taux, le smile se déplace dans l'autre sens, et que les couvertures par \omterm{def:dv:local-volatility:model}{volatilité locale} étaient alors pires que
l'absence de couverture~; le chapitre~11 construit le modèle qu'ils ont proposé.

\begin{omfigure}
\begin{tikzpicture}
\begin{axis}[width=10.5cm,height=5cm,
  xlabel={strike},ylabel={volatilité implicite à 3 mois (\%)},
  tick label style={font=\small},label style={font=\small},
  xmin=85,xmax=115,ymin=11,ymax=31,grid=major,grid style={gray!25},
  legend columns=2,legend style={font=\footnotesize,at={(0.5,-0.3)},anchor=north,draw=none,fill=none,/tikz/every even column/.append style={column sep=8pt}},
  table/col sep=comma]
\addplot[omDef,very thick,mark=*,mark size=1.3pt] table[x=strike,y=before]{figdata/derivatives/09-local-volatility/dynamics.csv};
\addplot[omThm,very thick,mark=square*,mark size=1.3pt] table[x=strike,y=after_lv]{figdata/derivatives/09-local-volatility/dynamics.csv};
\addplot[omProp,thick,dashed] table[x=strike,y=after_delta]{figdata/derivatives/09-local-volatility/dynamics.csv};
\legend{{avant, et après sous sticky strike},après une baisse de 5~\% sous volatilité locale,après sous sticky delta}
\end{axis}
\end{tikzpicture}

\omcaption{Le smile à trois mois avant et après une baisse de 5~\% du spot (de 100 à 95), selon trois règles.
La \omterm{def:dv:local-volatility:model}{volatilité locale} (carrés, par simulation sur nombres aléatoires communs) relève la volatilité de chaque strike~; \omterm{def:dv:implied-volatility-and-its-surface:sticky}{sticky
strike} laisse la courbe là où elle était~; \omterm{def:dv:implied-volatility-and-its-surface:sticky}{sticky delta} la déplace vers la gauche. Données~: le
tutoriel.}\label{fig:dv:local-volatility:dynamics}
\end{omfigure}

\begin{definition}[Smile forward]\label{def:dv:local-volatility:forward}
Le \emph{smile forward}\index{smile forward} de $t_1$ à $t_2$ est la volatilité implicite, en fonction de
le strike relatif $x$, des options à départ différé qui versent $(S_{t_2}-xS_{t_1})^+$ en $t_2$~: le smile que le
modèle s'attend à observer en $t_1$ pour des options de maturité $t_2-t_1$.
\end{definition}

Un modèle à dynamique homogène en temps (les modèles à volatilité stochastique du chapitre~10) produit des \omterm{def:dv:local-volatility:forward}{smiles forward}
qui ressemblent au smile du jour pour la même maturité. La \omterm{def:dv:local-volatility:model}{volatilité locale} ne le fait pas~: sa fonction de \omterm{def:dv:local-volatility:model}{volatilité locale} est
pentue près du spot du jour et aux temps courts, et s'aplatit à l'écart de ceux-ci, si bien que le smile qu'elle engendre pour une date future
est plus plat. Sur la surface du chapitre, le smile à six mois, vu six mois à l'avance, a un skew à la monnaie de
$-0.19$ par unité de \omterm{def:dv:implied-volatility-and-its-surface:surface}{log-moneyness}, contre $-0.34$ pour le smile à six mois du jour~: 57~\% de celui-ci
(la \cref{fig:dv:local-volatility:forward}). Son niveau est plus élevé (20.7~\% contre 17.8~\%), ce qui n'est que la
variance forward de la structure par terme. Les produits dont la valeur dépend des skews forward (options à départ différé,
cliquets, chapitre~16) sont donc valorisés différemment par des modèles qui s'accordent sur chaque vanille.

\begin{omfigure}
\begin{tikzpicture}
\begin{axis}[width=10.5cm,height=5cm,
  xlabel={strike relatif $x$},ylabel={volatilité implicite (\%)},
  tick label style={font=\small},label style={font=\small},
  xmin=0.8,xmax=1.2,ymin=12,ymax=28,grid=major,grid style={gray!25},
  legend columns=2,legend style={font=\footnotesize,at={(0.5,-0.3)},anchor=north,draw=none,fill=none,/tikz/every even column/.append style={column sep=8pt}},
  table/col sep=comma]
\addplot[omThm,very thick,mark=*,mark size=1.3pt] table[x=x,y=forward]{figdata/derivatives/09-local-volatility/forward_smile.csv};
\addplot[omDef,very thick] table[x=x,y=today]{figdata/derivatives/09-local-volatility/forward_smile.csv};
\legend{volatilité locale~: smile à 6 mois dans 6 mois,smile à 6 mois d'aujourd'hui}
\end{axis}
\end{tikzpicture}

\omcaption{Le \omterm{def:dv:local-volatility:forward}{smile forward} du \omterm{def:dv:local-volatility:model}{modèle de volatilité locale} (options de maturité six mois démarrant dans six mois, par
simulation), comparé au smile à six mois d'aujourd'hui. Le \omterm{def:dv:local-volatility:forward}{smile forward} est plus élevé, parce que la variance forward est
plus élevée, et plus plat~: son skew à la monnaie vaut 57~\% de celui d'aujourd'hui. Données~: le
tutoriel.}\label{fig:dv:local-volatility:forward}
\end{omfigure}

\section{Tutoriel~: d'une surface à un modèle de volatilité locale}\label{tut:dv:local-volatility:tut}

\begin{tutorial}
\textbf{Objectif.} Calculer la \omterm{def:dv:local-volatility:model}{volatilité locale} de Dupire à partir d'une surface SSVI, simuler le modèle, vérifier qu'il
reproduit les vanilles, et mesurer son \omterm{def:dv:local-volatility:forward}{smile forward} et la dynamique de son smile. \textbf{État final~:} les quatre figures
du chapitre et les chiffres du problème du week-end.
\begin{tutsteps}
\item \textbf{Variance locale} à partir de toute fonction régulière de variance totale~: la dérivée en temps divisée par le
facteur butterfly.
\omcode{code/firm/localvol/firm_localvol.py}{22}{32}{Variance locale de Dupire en variance implicite totale.}\label{lst:dv:local-volatility:dupire}
\item \textbf{Le générateur de trajectoires}~: un schéma log-Euler qui lit la \omterm{def:dv:local-volatility:model}{volatilité locale} au milieu de chaque
pas, fixé en spot absolu, afin qu'un mouvement du spot lise la fonction ailleurs.
\omcode{code/firm/localvol/firm_localvol.py}{61}{78}{Simulation du modèle de volatilité locale.}\label{lst:dv:local-volatility:sim}
\item \textbf{Exécuter} \texttt{dv\_localvol.mc\_smile()}, \texttt{forward\_smile()}, \texttt{dynamics()} et
\texttt{fig\_localvol.py}.
\end{tutsteps}
\textbf{Pour aller plus loin.} Construire la grille de \omterm{def:dv:local-volatility:model}{volatilité locale} en \omterm{def:dv:implied-volatility-and-its-surface:surface}{log-moneyness} par rapport au spot courant au lieu de
la référence forward de la calibration, et constater que le test de dynamique donne le mauvais signe (la première version du code de ce chapitre
l'a fait)~; puis diviser le pas de temps par deux et vérifier que les erreurs de repricing ne bougent pas.
\end{tutorial}

\section{Construction~: la surface de volatilité locale}\label{bld:dv:local-volatility:localvol}

\begin{build}
\bfield{Rôle} Le premier modèle de la firme miniature au-delà du \omterm{def:dv:black-scholes-three-ways:model}{modèle de Black--Scholes}~: calibré exactement sur la surface du
chapitre~8, utilisé pour valoriser les produits dépendant de la trajectoire (chapitres 15--18) et comme composante locale du modèle
à \omterm{def:dv:local-volatility:model}{volatilité locale} stochastique (chapitre~20).
\bfield{Interface} \texttt{local\_variance(w\_fn, k, t)}~; \texttt{build\_grid(w\_fn, times, ks, s\_ref, carry)
-> LocalVolGrid} avec \texttt{sigma(t, s)}~; \texttt{simulate(grid, s0, horizon, steps, n\_paths, seed,
record)}.
\bfield{Règles} Lever une erreur si le facteur butterfly est non positif ou si la dérivée en temps est négative~; grille en \omterm{def:dv:implied-volatility-and-its-surface:surface}{log-moneyness}
par rapport à la référence forward de la calibration (fixée en spot absolu)~; interpolation linéaire en temps et en moneyness, constante
au-delà~; paires antithétiques~; volatilité au milieu de chaque pas.
\bfield{Tests d'acceptation} \texttt{code/firm/localvol/tests/}~: une surface plate donne une \omterm{def:dv:local-volatility:model}{volatilité locale} plate égale
à elle-même~; une structure par terme sans skew donne les volatilités forward~; les vanilles à six mois simulées reproduisent la
surface à 0.12 point de volatilité près, et à 0.07 avec quatre pas par jour~; la grille lève une erreur sur une surface avec
un \omterm{def:dv:implied-volatility-and-its-surface:static}{arbitrage butterfly}.
\bfield{Pour aller plus loin} Un pricer par différences finies sur la grille de \omterm{def:dv:local-volatility:model}{volatilité locale} (chapitre~22), qui
supprime le bruit de Monte-Carlo de la vérification.
\end{build}

\begin{omsources}
\item B.~Dupire, ``Pricing with a smile'', \emph{Risk} 7(1) (1994) 18--20.
\item E.~Derman and I.~Kani, ``Riding on a smile'', \emph{Risk} 7(2) (1994) 32--39.
\item I.~Gyöngy, ``Mimicking the one-dimensional marginal distributions of processes having an Itô
differential'', \emph{Probability Theory and Related Fields} 71 (1986) 501--516.
\item P.~S.~Hagan, D.~Kumar, A.~S.~Lesniewski and D.~E.~Woodward, ``Managing smile risk'', \emph{Wilmott}
(2002).
\item L.~Bergomi, ``Smile dynamics'', \emph{Risk} (septembre 2004).
\item J.~Gatheral, \emph{The Volatility Surface}, Wiley, 2006, chapitres 1--3.
\end{omsources}

\section{Exercices}

\begin{exercise}[$\star$]\label{exo:dv:local-volatility:1}
En un point $(k,T)$ donné, la variance totale vaut $w=0.01$, sa dérivée en temps 0.045 et le facteur butterfly $g=0.9$.
Donner la \omterm{def:dv:local-volatility:model}{volatilité locale}.
\end{exercise}

\begin{exercise}[$\star$]\label{exo:dv:local-volatility:2}
La volatilité implicite à trois mois vaut 16.4~\% à la monnaie, avec une pente de $-0.47$ par unité de \omterm{def:dv:implied-volatility-and-its-surface:surface}{log-moneyness}. Utiliser
la règle de la \cref{prop:dv:local-volatility:twice} pour estimer la \omterm{def:dv:local-volatility:model}{volatilité locale} en $k=-0.1$, et la comparer
avec la valeur exacte de 27.4~\%.
\end{exercise}

\begin{exercise}[$\star$]\label{exo:dv:local-volatility:3}
Dans un modèle où la volatilité du spot vaut 15~\% ou 30~\% avec probabilité égale, indépendamment du spot, quelle est
sa \omterm{def:dv:local-volatility:projection}{projection markovienne}~? Et si, lorsque le spot vaut 90, la probabilité de l'état haut est 0.8~?
\end{exercise}

\begin{exercise}[$\star\star$]\label{exo:dv:local-volatility:4}
Expliquer pourquoi la surface implicite doit être exempte d'arbitrage avant d'appliquer la \omterm{thm:dv:local-volatility:dupire}{formule de Dupire}, et dire
ce qu'une cotation bruitée fait à la \omterm{def:dv:local-volatility:model}{volatilité locale} au voisinage de celle-ci.
\end{exercise}

\begin{exercise}[$\star\star$]\label{exo:dv:local-volatility:5}
À partir des variances totales à la monnaie à six mois et à un an sur la surface du chapitre (volatilités 17.79~\% et
19.19~\%), calculer la volatilité forward à la monnaie entre six mois et un an, et la comparer avec la
volatilité à départ différé à la monnaie du \omterm{def:dv:local-volatility:model}{modèle de volatilité locale}, 20.68~\%.
\end{exercise}

\begin{exercise}[$\star\star$]\label{exo:dv:local-volatility:6}
En utilisant la règle de la \cref{prop:dv:local-volatility:dynamics}, prévoir la variation de la volatilité à la monnaie à trois mois
après une baisse de 1~\%, et la comparer avec les 1.31 points simulés.
\end{exercise}

\begin{exercise}[$\star\star\star$]\label{exo:dv:local-volatility:7}
\emph{Programmation.} Repricer le put à six mois de strike 90 avec \texttt{mc\_smile} en utilisant 50\,000 et 200\,000 trajectoires~; comparer les
écarts aux erreurs types, et dire ce qui, outre le bruit de Monte-Carlo, peut produire un écart.
\end{exercise}

\begin{exercise}[$\star\star\star$]\label{exo:dv:local-volatility:8}
\emph{Trouver la faille.} «~Notre \omterm{def:dv:local-volatility:model}{modèle de volatilité locale} reproduit le prix de toute option cotée, donc il price tout produit dont le~
payoff dépend de l'indice à une date, y compris les options à départ différé.~»
\end{exercise}

\section{Problème~: Le smile forward qui s'aplatit}

\begin{problem}[{Problème du week-end --- ce que dit la volatilité locale sur le smile dans six mois}]\label{pb:dv:local-volatility:1}
Un structureur veut pricer un put à départ différé, de strike égal à 95~\% du niveau de l'indice dans six mois,
arrivant à échéance six mois plus tard, avec la surface du chapitre et des taux nuls. Le seul modèle calibré du desk est un \omterm{def:dv:local-volatility:model}{modèle de
volatilité locale}.

\medskip\noindent\textbf{Partie I --- La \omterm{def:dv:local-volatility:model}{volatilité locale}.}
\begin{enumerate}
\item Donner les volatilités implicite et locale à trois mois à la monnaie et en $k=-0.1$.
\item Donner le rapport de leurs pentes à la monnaie.
\item Pourquoi la \omterm{def:dv:local-volatility:model}{volatilité locale} est-elle plus pentue que la volatilité implicite~?
\item Donner les erreurs de repricing à six mois pour les strikes 80, 100 et 110.
\item Que signalerait une erreur de repricing d'un point de volatilité~?
\end{enumerate}

\medskip\noindent\textbf{Partie II --- Le \omterm{def:dv:local-volatility:forward}{smile forward}.}
\begin{enumerate}[resume]
\item Donner la volatilité à la monnaie de départ différé du modèle et la volatilité à la monnaie à six mois d'aujourd'hui.
\item Expliquer l'écart de niveau avec la variance forward.
\item Donner les skews à la monnaie du \omterm{def:dv:local-volatility:forward}{smile forward} et du smile à six mois d'aujourd'hui.
\item Donner leur rapport.
\item Pourquoi la \omterm{def:dv:local-volatility:model}{volatilité locale} aplatit-elle le \omterm{def:dv:local-volatility:forward}{smile forward}~?
\end{enumerate}

\medskip\noindent\textbf{Partie III --- Le produit.}
\begin{enumerate}[resume]
\item Donner la volatilité implicite du \omterm{def:dv:local-volatility:forward}{smile forward} en $x=0.95$, et la volatilité que donne en ce point la forme du smile
à six mois d'aujourd'hui, une fois décalée jusqu'au niveau à la monnaie forward.
\item Pricer le put à départ différé pour 100 de notionnel avec chacun des deux modèles (taux nuls, formule de Black sur un forward de 1).
\item Quel modèle un desk qui vend ce put préférerait-il, et pourquoi cela constitue-t-il un signal d'alerte~?
\item Quel instrument de marché permettrait de révéler quel skew forward est le bon~?
\item Comment constituer une réserve de modèle pour l'écart entre les modèles (chapitre 27)~?
\end{enumerate}

\medskip\noindent\textbf{Partie IV --- Jugement.}
\begin{enumerate}[resume]
\item Après une baisse de 5~\% de l'indice, que prédit la \omterm{def:dv:local-volatility:model}{volatilité locale} pour la volatilité à la monnaie à trois mois,
et que prédit le \omterm{def:dv:implied-volatility-and-its-surface:sticky}{sticky strike}~?
\item Quand la \omterm{def:dv:local-volatility:model}{volatilité locale} est-elle le bon modèle à utiliser~?
\item Que doit ajouter un modèle pour conserver des skews forward pentus~?
\item Énoncer le \emph{résultat nommé}~: le rapport entre le skew forward de la \omterm{def:dv:local-volatility:model}{volatilité locale} (à six mois dans six mois) et
le skew à six mois d'aujourd'hui.
\item En une phrase~: qu'est-ce que l'ajustement de toutes les options vanilles laisse indéterminé~?
\end{enumerate}
\end{problem}

\section{Questions d'entretien}

\begin{interviewq}[$\star$ \iqroles{chercheur}]\label{iq:dv:local-volatility:1}
Qu'est-ce que la \omterm{def:dv:local-volatility:model}{volatilité locale}, et pourquoi est-elle unique~?
\end{interviewq}

\begin{interviewq}[$\star$ \iqroles{trader}]\label{iq:dv:local-volatility:2}
L'indice chute. Que dit un \omterm{def:dv:local-volatility:model}{modèle de volatilité locale} de ce qui arrive à la volatilité implicite à la monnaie, et pourquoi~?
\end{interviewq}

\begin{interviewq}[$\star\star$ \iqroles{chercheur}]\label{iq:dv:local-volatility:3}
Dériver la \omterm{thm:dv:local-volatility:dupire}{formule de Dupire} avec des taux nuls.
\end{interviewq}

\begin{interviewq}[$\star\star$ \iqroles{chercheur, banque}]\label{iq:dv:local-volatility:4}
Deux modèles calibrent les mêmes options vanilles. Pour quels produits leurs prix peuvent-ils différer, et pour lesquels ne le peuvent-ils pas~?
\end{interviewq}

\begin{interviewq}[$\star\star$ \iqroles{développeur}]\label{iq:dv:local-volatility:5}
Votre \omterm{def:dv:implied-volatility-and-its-surface:surface}{surface de volatilité} locale présente des pics qui rendent instable le pricer Monte Carlo. D'où viennent-ils, et
comment les corrigez-vous~?
\end{interviewq}

\begin{interviewq}[$\star\star\star$ \iqroles{chercheur, trader}]\label{iq:dv:local-volatility:6}
Qu'est-ce qu'une \omterm{def:dv:local-volatility:projection}{projection markovienne}, et comment l'utiliseriez-vous pour calibrer un modèle à volatilité stochastique
sur l'ensemble de la surface~?
\end{interviewq}
